% Общий исходник: условия + решения.
% Подключается из 2026.08.28_imc_tasks.tex (\sol=0) и 2026.08.28_imc_solutions.tex (\sol=1).

\providecommand{\src}[1]{\ \textup{\small\textit{(IMC #1)}}}
\providecommand{\Zh}{\mathbb{Z}}
\providecommand{\Rh}{\mathbb{R}}

%%%%%%%%%%%%%%%%%%%%%%%%%%%%%%%%%%%%%%%%%%%%%%%%%%%%%%%%%%%%%%%%%%%%%%%%%%%%%%
\q Пусть $n$ --- натуральное число. Сколько существует слов длины $n$ в алфавите $\{a,b,c,d\}$, в которых буква $a$ встречается чётное число раз и буква $b$ тоже встречается чётное число раз? (Например, при $n=2$ таких слов шесть: $aa$, $bb$, $cc$, $dd$, $cd$, $dc$.)
\src{2020, день 1, задача 1}

\if\sol1
\textbf{Решение.} \emph{Ответ:} $4^{n-1}+2^{n-1}$.

Для слова $w$ обозначим через $\alpha(w)$ и $\beta(w)$ количества букв $a$ и $b$ в нём и положим
\[
\varepsilon(w)=(-1)^{\alpha(w)},\qquad \delta(w)=(-1)^{\beta(w)}.
\]
Выражение $\tfrac14\bigl(1+\varepsilon(w)\bigr)\bigl(1+\delta(w)\bigr)$ равно $1$, если оба числа $\alpha(w)$ и $\beta(w)$ чётны, и равно $0$ иначе. Поэтому искомое количество равно
\[
\frac14\sum_{w}\bigl(1+\varepsilon(w)+\delta(w)+\varepsilon(w)\delta(w)\bigr),
\]
где сумма берётся по всем $4^n$ словам. Каждая из четырёх сумм считается по буквам независимо (произведение по позициям):
\[
\sum_w 1=4^n,\qquad
\sum_w\varepsilon(w)=\bigl((-1)+1+1+1\bigr)^n=2^n,\qquad
\sum_w\delta(w)=2^n,
\]
а для $\varepsilon\delta$ вклады букв $a,b,c,d$ равны $-1,-1,+1,+1$, их сумма равна нулю, поэтому $\sum_w\varepsilon(w)\delta(w)=0^n=0$. Итого получаем
\[
\frac{4^n+2^n+2^n+0}{4}=4^{n-1}+2^{n-1}.
\]
При $n=2$: $4+2=6$ --- совпадает с примером.
\fi

%%%%%%%%%%%%%%%%%%%%%%%%%%%%%%%%%%%%%%%%%%%%%%%%%%%%%%%%%%%%%%%%%%%%%%%%%%%%%%
\q Для натурального $n$ обозначим через $p(n)$ число способов представить $n$ в виде суммы натуральных слагаемых (порядок слагаемых не важен). Например, $p(4)=5$, поскольку
\[
4=3+1=2+2=2+1+1=1+1+1+1 .
\]
Положим также $p(0)=1$. Докажите, что $p(n)-p(n-1)$ равно числу представлений числа $n$ в виде суммы слагаемых, каждое из которых строго больше $1$.
\src{2012, день 1, задача 1}

\if\sol1
\textbf{Решение.} Все разбиения числа $n$ делятся на два класса: содержащие слагаемое $1$ и не содержащие его. Второй класс --- это в точности разбиения на слагаемые, большие $1$, и надо доказать, что их ровно $p(n)-p(n-1)$. Значит, достаточно проверить, что разбиений первого класса ровно $p(n-1)$.

Построим биекцию между разбиениями числа $n-1$ и разбиениями числа $n$, содержащими единицу: разбиению $n-1=a_1+\dots+a_k$ сопоставим разбиение $n=a_1+\dots+a_k+1$. Это действительно разбиение числа $n$ с единицей; обратное отображение --- выбросить одну единицу --- определено однозначно. Следовательно, классов первого типа ровно $p(n-1)$, и
\[
\#\{\text{разбиения без единиц}\}=p(n)-p(n-1). \qquad \blacksquare
\]
(При $n=1$ равенство тоже верно: $p(1)-p(0)=0$, а сумм из слагаемых $>1$, дающих $1$, нет.)
\fi

%%%%%%%%%%%%%%%%%%%%%%%%%%%%%%%%%%%%%%%%%%%%%%%%%%%%%%%%%%%%%%%%%%%%%%%%%%%%%%
\q Иван написал на доске матрицу $\begin{pmatrix} 2 & 3\\ 2 & 4\end{pmatrix}$. Затем он несколько раз проделывает такую операцию: выбирает строку или столбец матрицы и покомпонентно умножает или делит её (его) на другую строку или, соответственно, на другой столбец. Может ли Иван за конечное число ходов получить матрицу $\begin{pmatrix} 2 & 4\\ 2 & 3\end{pmatrix}$?
\src{2023, день 2, задача 6}

\if\sol1
\textbf{Решение.} \emph{Ответ:} нет.

Все элементы исходной матрицы положительны, а разрешённые операции сохраняют положительность, поэтому логарифмы элементов определены на всех шагах. Сопоставим матрице $X=(x_{ij})$ матрицу логарифмов
\[
L(X)=\bigl(\log_2 x_{ij}\bigr).
\]
Покомпонентное умножение строки на другую строку превращается в прибавление одной строки матрицы $L$ к другой, деление --- в вычитание; то же самое со столбцами. Но элементарные преобразования «прибавить к одной строке другую» и «прибавить к одному столбцу другой» не меняют определителя. Значит, величина $\det L(X)$ --- инвариант игры.

Для начальной матрицы $A$ и целевой матрицы $B$:
\[
\det L(A)=1\cdot 2-\log_2 3\cdot 1=2-\log_2 3=\log_2\tfrac43>0,
\]
\[
\det L(B)=1\cdot\log_2 3-2\cdot 1=\log_2 3-2=-\log_2\tfrac43<0 .
\]
Инварианты различны, поэтому получить $B$ из $A$ невозможно. \hfill$\blacksquare$
\fi

%%%%%%%%%%%%%%%%%%%%%%%%%%%%%%%%%%%%%%%%%%%%%%%%%%%%%%%%%%%%%%%%%%%%%%%%%%%%%%
\q Найдите все пары комплексных чисел $(a,b)$, для которых $|a|=|b|=1$ и $a+b+a\bar b\in\Rh$.
\src{2024, день 1, задача 1}

\if\sol1
\textbf{Решение.} \emph{Ответ:} $a=1$ (при любом $b$), либо $b=-1$ (при любом $a$), либо $b=-a$.

Так как $|a|=|b|=1$, то $\bar a=1/a$ и $\bar b=1/b$, а $a\bar b=a/b$. Комплексное число вещественно тогда и только тогда, когда оно совпадает со своим сопряжением, поэтому условие равносильно
\[
a+b+\frac ab=\frac1a+\frac1b+\frac ba .
\]
Умножим обе части на $ab\ne 0$:
\[
a^2b+ab^2+a^2=b+a+b^2,
\qquad\text{то есть}\qquad
a^2b+ab^2+a^2-a-b-b^2=0 .
\]
Левая часть раскладывается на множители:
\[
a^2(b+1)+a(b^2-1)-b(b+1)=(b+1)\bigl(a^2+a(b-1)-b\bigr)=(b+1)(a+b)(a-1).
\]
Итак, условие равносильно $(b+1)(a+b)(a-1)=0$, то есть $b=-1$, или $b=-a$, или $a=1$.

Проверка: при $a=1$ получаем $1+b+\bar b=1+2\operatorname{Re}b\in\Rh$; при $b=-1$ получаем $a-1-a=-1\in\Rh$; при $b=-a$ получаем $a-a-a\bar a=-|a|^2=-1\in\Rh$. \hfill$\blacksquare$
\fi

%%%%%%%%%%%%%%%%%%%%%%%%%%%%%%%%%%%%%%%%%%%%%%%%%%%%%%%%%%%%%%%%%%%%%%%%%%%%%%
\q Все стороны и диагонали правильного $43$-угольника покрашены в красный или синий цвет так, что из каждой вершины выходит ровно $20$ красных и ровно $22$ синих отрезка. Треугольник с вершинами в вершинах многоугольника называется \emph{одноцветным}, если все три его стороны одного цвета. Известно, что синих одноцветных треугольников ровно $2022$. Сколько красных одноцветных треугольников?
\src{2022, день 2, задача 5}

\if\sol1
\textbf{Решение.} \emph{Ответ:} $859$.

Назовём \emph{разноцветным углом} пару отрезков разного цвета, выходящих из одной вершины. Посчитаем число разноцветных углов двумя способами.

С одной стороны, из каждой вершины выходит $20$ красных и $22$ синих отрезка, поэтому в каждой вершине ровно $20\cdot 22=440$ разноцветных углов, а всего их
\[
43\cdot 440=18\,920 .
\]

С другой стороны, посмотрим на треугольник. Если он одноцветный, разноцветных углов в нём нет. Если же он не одноцветный, то среди трёх его сторон есть стороны обоих цветов, скажем, две одного цвета и одна другого; тогда ровно две его вершины являются концами разноцветной пары сторон. Значит, каждый неодноцветный треугольник даёт ровно $2$ разноцветных угла, и число неодноцветных треугольников равно
\[
\frac{18\,920}{2}=9460 .
\]

Всего треугольников $\binom{43}{3}=\frac{43\cdot 42\cdot 41}{6}=12\,341$, поэтому одноцветных треугольников
\[
12\,341-9460=2881 .
\]
Из них синих $2022$, значит, красных
\[
2881-2022=859 . \qquad \blacksquare
\]
\fi

%%%%%%%%%%%%%%%%%%%%%%%%%%%%%%%%%%%%%%%%%%%%%%%%%%%%%%%%%%%%%%%%%%%%%%%%%%%%%%
\q Пусть $z$ --- комплексное число, для которого $|z+1|>2$. Докажите, что $|z^3+1|>1$.
\src{2013, день 2, задача 1}

\if\sol1
\textbf{Решение.} Так как $z^3+1=(z+1)(z^2-z+1)$ и $|z+1|>2$, достаточно доказать, что
\[
|z^2-z+1|>\tfrac12 .
\]
Обозначим $w=z+1$, тогда $z=w-1$ и
\[
z^2-z+1=(w-1)^2-(w-1)+1=w^2-3w+3 .
\]
Пусть $r=|w|>2$ и $t=\operatorname{Re}w$, так что $w+\bar w=2t$ и $|t|\le r$. Тогда
\[
w^2+\bar w^2=(w+\bar w)^2-2|w|^2=4t^2-2r^2,
\]
и, раскрывая произведение $(w^2-3w+3)(\bar w^2-3\bar w+3)$, получаем
\[
|w^2-3w+3|^2=r^4-3r^2(w+\bar w)+3(w^2+\bar w^2)+9r^2-9(w+\bar w)+9
\]
\[
=r^4-6r^2t+3(4t^2-2r^2)+9r^2-18t+9
=12t^2-(6r^2+18)\,t+\bigl(r^4+3r^2+9\bigr).
\]
Это квадратный трёхчлен относительно $t$ с положительным старшим коэффициентом; его минимум достигается при $t=\frac{r^2+3}{4}$ и равен
\[
r^4+3r^2+9-\frac{(6r^2+18)^2}{48}=r^4+3r^2+9-\frac34\bigl(r^2+3\bigr)^2=\frac{(r^2-3)^2}{4}.
\]
Следовательно,
\[
|z^2-z+1|^2=|w^2-3w+3|^2\ \ge\ \frac{(r^2-3)^2}{4}.
\]
Так как $r>2$, то $r^2-3>1$, значит, $|z^2-z+1|^2>\tfrac14$ и $|z^2-z+1|>\tfrac12$. Окончательно
\[
|z^3+1|=|z+1|\cdot|z^2-z+1|>2\cdot\tfrac12=1 . \qquad \blacksquare
\]
\fi

%%%%%%%%%%%%%%%%%%%%%%%%%%%%%%%%%%%%%%%%%%%%%%%%%%%%%%%%%%%%%%%%%%%%%%%%%%%%%%
\q В школе учатся $2n$ учеников ($n\ge 2$). Каждую неделю ровно $n$ учеников едут на экскурсию. Через несколько недель оказалось, что любые два ученика хотя бы раз побывали на экскурсии вместе. Какое наименьшее число экскурсий для этого необходимо?
\src{2013, день 1, задача 3}

\if\sol1
\textbf{Решение.} \emph{Ответ:} $6$ (при любом $n\ge2$).

\emph{Оценка.} На одной экскурсии ученик встречается ровно с $n-1$ другими, а познакомиться ему нужно со всеми $2n-1$ одноклассниками. Поскольку $2(n-1)=2n-2<2n-1$, двух экскурсий каждому ученику не хватает, то есть каждый ездил хотя бы трижды. Просуммируем по всем ученикам: общее число «участий» не меньше $3\cdot 2n=6n$. С другой стороны, на каждой экскурсии ровно $n$ участников, поэтому если экскурсий было $m$, то число участий равно $mn$. Из $mn\ge 6n$ получаем $m\ge 6$.

\emph{Пример на $6$ экскурсий.} Любое $n\ge2$ представимо в виде $n=2x+3y$ с целыми $x,y\ge0$ (при чётном $n$ берём $y=0$, при нечётном $n\ge3$ берём $y=1$ и $x=\frac{n-3}{2}$). Разобьём учеников на четыре группы $A,B,C,D$ по $x$ человек и шесть групп $E,F,G,H,I,J$ по $y$ человек: всего $4x+6y=2n$ человек. Проведём экскурсии
\[
ABEFG,\quad CDEFH,\quad ACGHI,\quad BDGHJ,\quad ADEIJ,\quad BCFIJ,
\]
где запись означает объединение указанных групп. В каждой экскурсии ровно $2x+3y=n$ человек.

Проверим, что любые два ученика встретились. Группы $A,B,C,D$ входят в экскурсии парами $AB$, $CD$, $AC$, $BD$, $AD$, $BC$ --- это все шесть пар, а каждая из этих групп участвует в трёх экскурсиях, поэтому встретились и ученики внутри одной группы. Группы $E,\dots,J$ входят тройками $EFG$, $EFH$, $GHI$, $GHJ$, $EIJ$, $FIJ$; несложная проверка показывает, что все $15$ пар этих групп встречаются, и каждая группа участвует в трёх экскурсиях. Наконец, каждая из групп $A,B,C,D$ встречается со всеми шестью группами $E,\dots,J$: например, $A$ ездит в экскурсиях $ABEFG$, $ACGHI$, $ADEIJ$, и в них участвуют все группы $E,F,G,H,I,J$; для $B$, $C$, $D$ проверка такая же. \hfill$\blacksquare$
\fi

%%%%%%%%%%%%%%%%%%%%%%%%%%%%%%%%%%%%%%%%%%%%%%%%%%%%%%%%%%%%%%%%%%%%%%%%%%%%%%
\q Для натурального $n$ пусть $f(n)$ --- число, которое получается, если записать $n$ в двоичной системе и заменить каждый нуль единицей, а каждую единицу нулём. Например, $23=10111_2$, после замены получаем $01000_2=1000_2$, поэтому $f(23)=8$. Докажите, что
\[
\sum_{k=1}^{n} f(k)\ \le\ \frac{n^2}{4},
\]
и выясните, при каких $n$ достигается равенство.
\src{2015, день 1, задача 2}

\if\sol1
\textbf{Решение.} Ключевое наблюдение: если $2^{r-1}\le k<2^{r}$, то в двоичной записи числа $k$ ровно $r$ цифр, и при замене цифр каждый разряд пары $(k,f(k))$ даёт единицу, то есть
\[
k+f(k)=\underbrace{11\ldots1}_{r\ \text{единиц}}\ (\text{в двоичной записи})\ =\ 2^r-1 .
\]

Пусть $n$ имеет $s$ двоичных цифр, то есть $2^{s-1}\le n\le 2^{s}-1$. Разобьём отрезок $1\le k\le n$ на блоки по числу цифр:
\[
\sum_{k=1}^{n}\bigl(k+f(k)\bigr)
=\sum_{r=1}^{s-1}\ \sum_{2^{r-1}\le k<2^{r}}(2^r-1)\ +\sum_{2^{s-1}\le k\le n}(2^s-1)
\]
\[
=\sum_{r=1}^{s-1}2^{r-1}\bigl(2^r-1\bigr)+\bigl(n-2^{s-1}+1\bigr)\bigl(2^s-1\bigr).
\]
Обозначим $q=2^s$. Первая сумма равна
\[
\sum_{r=1}^{s-1}\bigl(2^{2r-1}-2^{r-1}\bigr)=\frac{4^{s}-4}{6}-\bigl(2^{s-1}-1\bigr)=\frac{q^2-4}{6}-\frac q2+1 .
\]
Вычитая $\sum_{k=1}^n k=\frac{n(n+1)}2$ и приводя подобные, после несложных преобразований получаем
\[
\frac{n^2}{4}-\sum_{k=1}^{n}f(k)
=\frac34 n^2-\Bigl(q-\frac32\Bigr)n+\frac{q^2}{3}-q+\frac23
=\frac34\left(n-\frac{2q-2}{3}\right)\left(n-\frac{2q-4}{3}\right).
\]

Оба корня $\frac{2q-2}{3}=\frac{2^{s+1}-2}{3}$ и $\frac{2q-4}{3}=\frac{2^{s+1}-4}{3}$ отличаются друг от друга на $\frac23<1$, причём ровно один из них --- целое число: при чётном $s$ имеем $2^{s+1}\equiv2\pmod 3$ и целым является первый, при нечётном $s$ --- второй. Поэтому для целого $n$ два множителя либо оба одного знака (и произведение положительно), либо один из них равен нулю. В любом случае правая часть неотрицательна, что и даёт требуемое неравенство.

Равенство достигается ровно тогда, когда $n$ совпадает с целым корнем, то есть
\[
n=\frac{2^{s+1}-2}{3}\ \ (s\ \text{чётно})\qquad\text{или}\qquad n=\frac{2^{s+1}-4}{3}\ \ (s\ \text{нечётно}),
\]
то есть при $n=2,\ 4,\ 10,\ 20,\ 42,\ 84,\ \ldots$ \hfill$\blacksquare$

\emph{Проверка на маленьких числах:} $f(1)=0$, $f(2)=1$, $f(3)=0$, $f(4)=3$, $f(5)=2$. При $n=2$: $\sum f=1=\frac{2^2}{4}$; при $n=4$: $\sum f=4=\frac{4^2}{4}$; при $n=3$: $\sum f=1<\frac94$.
\fi

%%%%%%%%%%%%%%%%%%%%%%%%%%%%%%%%%%%%%%%%%%%%%%%%%%%%%%%%%%%%%%%%%%%%%%%%%%%%%%
\q Число $n$ называется \emph{совершенным}, если $\sigma(n)=2n$, где $\sigma(n)$ --- сумма всех натуральных делителей числа $n$. Пусть $n>6$ --- совершенное число и
\[
n=p_1^{e_1}p_2^{e_2}\cdots p_k^{e_k},\qquad p_1<p_2<\dots<p_k,
\]
--- его каноническое разложение на простые множители. Докажите, что показатель $e_1$ чётен.
\src{2014, день 1, задача 4}

\if\sol1
\textbf{Решение.} Предположим противное: пусть $e_1$ нечётно. Как известно,
\[
\sigma(n)=\prod_{i=1}^{k}\bigl(1+p_i+p_i^2+\dots+p_i^{e_i}\bigr)=2n .
\]
Так как $e_1$ нечётно, в первом сомножителе чётное число слагаемых ($e_1+1$ штук), и их можно разбить на пары:
\[
1+p_1+\dots+p_1^{e_1}=(1+p_1)\bigl(1+p_1^2+p_1^4+\dots+p_1^{e_1-1}\bigr).
\]
Значит, $p_1+1$ делит $\sigma(n)=2n=2p_1^{e_1}\cdots p_k^{e_k}$.

Покажем, что у числа $p_1+1\ge3$ есть нечётный простой делитель. Если $p_1=2$, то $p_1+1=3$. Если же $p_1$ нечётно, то $p_1$ --- наименьший простой делитель $n$, поэтому $n$ нечётно и $2n$ делится на $2$ ровно в первой степени; если бы $p_1+1$ было степенью двойки, из $(p_1+1)\mid 2n$ следовало бы $p_1+1\le2$, что невозможно.

Итак, пусть $q$ --- нечётный простой делитель числа $p_1+1$. Тогда $q\mid 2n$ и $q$ нечётно, поэтому $q\in\{p_1,\dots,p_k\}$. При этом $q\ne p_1$ (число $p_1$ не делит $p_1+1$), а для $i\ge3$ имеем $p_i>p_2>p_1$, откуда $p_i\ge p_1+2>p_1+1\ge q$, то есть $q\ne p_i$. Остаётся $q=p_2$. Но $p_2>p_1$ даёт $p_2\ge p_1+1$, а из $p_2\mid p_1+1$ следует $p_2\le p_1+1$. Значит, $p_2=p_1+1$: два последовательных натуральных числа простые, поэтому $p_1=2$, $p_2=3$ и $6\mid n$.

Теперь числа
\[
1,\quad \frac n6,\quad \frac n3,\quad \frac n2,\quad n
\]
--- пять попарно различных делителей числа $n$ (различны, так как $n>6$ и, значит, $\frac n6>1$). Поэтому
\[
\sigma(n)\ \ge\ 1+\frac n6+\frac n3+\frac n2+n=1+2n>2n,
\]
что противоречит совершенности числа $n$. Полученное противоречие доказывает, что $e_1$ чётно. \hfill$\blacksquare$
\fi

%%%%%%%%%%%%%%%%%%%%%%%%%%%%%%%%%%%%%%%%%%%%%%%%%%%%%%%%%%%%%%%%%%%%%%%%%%%%%%
\q Найдите все простые числа $p$, для которых существует ровно одно $a\in\{1,2,\dots,p\}$ такое, что число $a^3-3a+1$ делится на $p$.
\src{2020, день 2, задача 6}

\if\sol1
\textbf{Решение.} \emph{Ответ:} $p=3$.

Обозначим $f(x)=x^3-3x+1$ и будем работать с вычетами по модулю $p$. Пусть $a$ --- корень, то есть $p\mid f(a)$. Разделим $f$ на $x-a$:
\[
f(x)=(x-a)\,g(x)+f(a),\qquad g(x)=x^2+ax+\bigl(a^2-3\bigr),
\]
то есть по модулю $p$ имеем $f(x)\equiv(x-a)g(x)$.

\emph{Ключевое наблюдение:} число $b=a^2-2$ также является корнем $f$ по модулю $p$. Действительно, прямое вычисление даёт
\[
g\bigl(a^2-2\bigr)=\bigl(a^2-2\bigr)^2+a\bigl(a^2-2\bigr)+\bigl(a^2-3\bigr)
=a^4+a^3-3a^2-2a+1=(a+1)\,f(a),
\]
а правая часть делится на $p$. (Это отражение тождества $2\cos 2t=(2\cos t)^2-2$: корни многочлена $x^3-3x+1$ равны $2\cos\frac{2\pi}{9},\,2\cos\frac{4\pi}{9},\,2\cos\frac{8\pi}{9}$, и отображение $x\mapsto x^2-2$ переставляет их по циклу.)

По теореме Виета для $g$ третий корень равен $c=-a-b=-a^2-a+2$.

Пусть теперь корень единствен. Тогда все три корня совпадают, в частности $a\equiv b=a^2-2$, то есть
\[
a^2-a-2\equiv0,\qquad (a-2)(a+1)\equiv0 \pmod p,
\]
откуда $a\equiv2$ или $a\equiv-1$. Но
\[
f(2)=8-6+1=3,\qquad f(-1)=-1+3+1=3,
\]
и в обоих случаях $0\equiv f(a)=3\pmod p$, то есть $p=3$.

Осталось проверить, что $p=3$ подходит: $f(1)=-1$, $f(2)=3$, $f(3)=19$, и из этих трёх чисел на $3$ делится только $f(2)$. Значит, для $p=3$ подходящее $a$ существует и единственно. \hfill$\blacksquare$
\fi
